\documentclass[preprint,12pt]{elsarticle}
\usepackage[utf8]{inputenc}
\usepackage[T1]{fontenc}
\usepackage{amssymb}
\usepackage{amsmath}
\usepackage{float, graphicx, caption, xcolor, comment}
\usepackage{subcaption}
\captionsetup[sub]{labelformat=parens, labelsep=none, font=bf, position=top}
\usepackage{booktabs}
\usepackage{textcomp}
\usepackage{soul, url}
%\usepackage{lineno}
\usepackage{multirow}
\usepackage[margin=1in]{geometry}
\usepackage{hyperref, lineno}

\journal{Journal of Environmental Radioactivity}

%% ── Graphic path ──────────────────────────────────────────
\graphicspath{{./figure/}}

\begin{document}

\begin{frontmatter}

%% ── Title ─────────────────────────────────────────────────
\title{Low-Methoxyl Pectin--Ca$^{2+}$ Gelation as a Sample
Solidification Protocol for Marine Radioactivity Monitoring
by HPGe Gamma Spectrometry}

%% ── Authors ───────────────────────────────────────────────
\author[inst1]{Yejin Lee\fnref{fn1}}   %%jinfriend98@naver.com
\author[inst2]{Hyunsung Kim\fnref{fn1}}   %%khs030416@snu.ac.kr
\author[inst1]{Myeongsik Shin}
\author[inst1]{Gyoung Jun Choi}   %%cgj0212@hanmail.net
\author[inst1]{Jun Yoo}
\author[inst1]{Ku Kang\corref{cor1}}

\affiliation[inst1]{organization={CBRN Defense Research Institute},
            city={Seoul 06796},
            country={Republic of Korea}}

\affiliation[inst2]{organization={Department of Nuclear Engineering,
            Seoul National University},
            city={Seoul 08826},
            country={Republic of Korea}}

\cortext[cor1]{Corresponding author.
  E-mail: \href{mailto:bisu9082@gmail.com}{bisu9082@gmail.com}}

\fntext[fn1]{These authors contributed equally to this work.}

%% ── Abstract ──────────────────────────────────────────────
\begin{abstract}
Marine radioactivity monitoring requires rapid, reliable sample preparation
methods that remain functional aboard ships and other mobile platforms.
Liquid seawater samples, however, are prone to geometric instability and
spillage during vessel motion, which degrade HPGe gamma spectrometry
measurement quality.
We developed a solidification procedure based on low-methoxyl (LM) pectin
and controlled Ca$^{2+}$ supplementation, which converts seawater into a
mechanically stable, leak-resistant gel through egg-box cross-linking.
Under the optimised condition (1.4~g LM-pectin per 28.26~g seawater,
300~rpm, 120~min), the gel fills a cylindrical HPGe counting geometry
(diameter 60~mm, fill height 10~mm) reproducibly.
Paired recovery rates, calculated from measurements before and after
solidification of Korea Institute of Nuclear Safety (KINS)
proficiency-test samples, were
99.9~$\pm$~5.6\% for $^{137}$Cs ($n=5$),
88.8~$\pm$~12.6\% for $^{134}$Cs ($n=4$),
93.7~$\pm$~5.6\% for $^{241}$Am ($n=5$),
and 94.8~$\pm$~2.7\% for $^{60}$Co ($n=3$)
across a gamma-energy span of 59.5--1173.2~keV.
Sixteen of 17 paired measurements remained within 80--120\% recovery;
the exception was one $^{134}$Cs result (71.0\%).
The procedure uses food-grade reagents only, requires no pre-evaporation
or specialised equipment, and completes gelation within two hours,
making it a practical candidate for further validation in shipboard
environmental radioactivity surveys.
\end{abstract}

\begin{keyword}
Marine radioactivity monitoring \sep
Sample solidification \sep
Low-methoxyl pectin \sep
Egg-box gelation \sep
HPGe gamma spectrometry \sep
Field-deployable measurement
\end{keyword}

\end{frontmatter}
\linenumbers

%% ══════════════════════════════════════════════════════════
\section{Introduction}
%% ══════════════════════════════════════════════════════════

The Fukushima Daiichi nuclear accident in 2011 released substantial quantities
of radionuclides into the Pacific Ocean, demonstrating that rapid, on-site
seawater radioactivity measurement is not merely a routine task but a genuine
operational necessity~\cite{Buesseler2012, Lu2021}.
In the years since, the need to characterise marine radioactivity at sea---
rather than after samples reach a shore laboratory---has grown considerably,
driven by post-accident environmental surveillance programmes,
naval radiological protection, and international monitoring obligations
under IAEA frameworks~\cite{IAEA2005, Abbasi2021, Uddin2020}.
Long-term environmental radionuclide datasets also show the value of
measurement workflows that can be deployed consistently across monitoring
locations and conditions~\cite{Bae2025}.
High-purity germanium (HPGe) gamma spectrometry is the technique of choice
for simultaneous identification and quantification of multiple gamma-emitting
radionuclides in seawater without chemical separation, owing to its energy
resolution and broad nuclide coverage~\cite{Jerome2014, Debertin1988}.

The conventional approach aboard research or naval vessels is to collect
seawater into a cylindrical counting container and place the container
directly on the HPGe detector for measurement.
This works well in a shore laboratory, where the detector and sample remain
still.
At sea, however, persistent vessel motion causes the liquid surface to oscillate,
continuously changing the distance and solid angle between the sample and the
detector crystal (Fig.~\ref{fig:concept})~\cite{Corbacho2022, Choi2024}.
Even moderate swell introduces geometry-dependent efficiency errors of several
percent~\cite{Barrera2017}, making quantitative measurement
unreliable without lengthy post-measurement corrections that are impractical
in an emergency response context.
Beyond measurement accuracy, a liquid sample also poses a spillage risk:
in a confined shipboard instrument room, any leak of seawater containing
elevated radionuclide concentrations represents both a contamination hazard
and an operational problem~\cite{Lu2021, Choi2023}.

\begin{figure}[ht]
  \centering
  \includegraphics[width=0.6\linewidth]{Fig1.png}
  \caption{Schematic illustration of the sample solidification concept
    for field-deployable radionuclide measurement.
    (a) Conventional liquid seawater samples on mobile platforms
    (e.g., vessels) suffer from geometric instability and leakage risk,
    which compromise HPGe detector counting efficiency.
    (b) Pectin--Ca$^{2+}$ gelation solidifies the sample via an
    egg-box cross-linking mechanism
    ($[\text{C}_6\text{H}_8\text{O}_6]_n + \text{Ca}^{2+}
    \rightarrow (\text{Pectin--Ca--Pectin})_n$),
    limiting geometric variation and leakage risk while retaining
    radionuclides in the measurement sample.}
  \label{fig:concept}
\end{figure}

The logical solution is to convert the liquid sample into a solid or
gel-like form before measurement, so that the counting geometry is fixed
regardless of vessel motion.
Cement, epoxy resins~\cite{Costa2021}, and polymer
hydrogels~\cite{Du2020} have each been considered as solidification matrices
for radioactive materials; however, these methods share practical drawbacks
that make them poorly suited to shipboard use---elevated processing temperatures,
toxic or irritating reagents, curing times exceeding 24~h, or
difficulties achieving a spatially uniform distribution of dissolved
radionuclides~\cite{Rahman2007}.
What is needed is a solidification agent that works at room temperature,
uses safe reagents, and gels within a few hours.

We identified low-methoxyl (LM) pectin as a candidate because seawater
already contains Ca$^{2+}$, the principal cross-linking cation, and because
the gelation condition can be standardised by a small, controlled addition
of CaCl$_2\cdot$2H$_2$O.
LM-pectin (degree of esterification $<$50\%) gels in the presence of
divalent cations through the egg-box mechanism, in which Ca$^{2+}$ ions
coordinate with carboxylate groups on adjacent pectin chains, forming a
three-dimensional network:
$[\text{C}_6\text{H}_8\text{O}_6]_n + \text{Ca}^{2+}
\rightarrow (\text{Pectin--Ca--Pectin})_n$~\cite{Braccini2001, Grant1973}.
Seawater naturally contains approximately 400~mg~L$^{-1}$ Ca$^{2+}$ and
1300~mg~L$^{-1}$ Mg$^{2+}$~\cite{Millero2008}.
In this study, supplemental CaCl$_2\cdot$2H$_2$O was added to reduce
condition-to-condition variation and obtain a reproducible preparation
protocol.
LM-pectin is a food-grade polysaccharide widely used in the food
industry~\cite{May1990}, non-toxic, and stable at room temperature,
all properties that matter in a shipboard setting where chemical safety
and simplicity are priorities.

To our knowledge, LM-pectin--Ca$^{2+}$ gelation has not previously been
applied to seawater sample preparation for HPGe gamma spectrometry.
Open questions include how quantitatively the method preserves the activity
of dissolved radionuclides ($^{137}$Cs, $^{134}$Cs, $^{241}$Am, $^{60}$Co)
across the full energy range relevant to marine monitoring (59.5--1173.2~keV),
and whether the gel matrix introduces energy-dependent self-absorption
for low-energy emitters such as $^{241}$Am.
This study addresses those questions by: (i)~optimising pectin concentration
and gelation conditions for artificial seawater in a standard cylindrical
HPGe counting geometry; (ii)~quantifying radionuclide recovery and
measurement accuracy against Korea Institute of Nuclear Safety (KINS)
certified proficiency-test reference standards; and (iii)~assessing
analytical reproducibility across multiple sample preparations.
The result is the first validated protocol for pectin-gel solidification
of seawater samples intended for field HPGe gamma spectrometry.


%% ══════════════════════════════════════════════════════════
\section{Materials and Methods}
%% ══════════════════════════════════════════════════════════

\subsection{Reagents and Materials}

LM-pectin powder (degree of esterification $<$50\%, food grade;
105 RAPID SET, Korea) was used as received without further purification.
Supplemental Ca$^{2+}$ was added as calcium chloride dihydrate
(CaCl$_2\cdot$2H$_2$O, $\geq$99\%; DAEJUNG, Korea).
The added amount (0.0323~g per 28.26~g sample) corresponds to
approximately 312~mg~L$^{-1}$ supplemental Ca$^{2+}$.
Artificial seawater was prepared in-house following ASTM D1141 at
salinity 35~psu; its Ca$^{2+}$ concentration was approximately
400~mg~L$^{-1}$, consistent with the reference composition of standard
seawater~\cite{Millero2008}.
Radionuclide reference solutions ($^{137}$Cs, $^{134}$Cs, $^{241}$Am,
$^{60}$Co) were supplied by the Korea Institute of Nuclear Safety (KINS)
through their proficiency testing programme.
We used samples from three rounds (2019, 2022, 2024) spanning a 1.8-fold
range in reference activity concentration (24.6--45.1~Bq~kg$^{-1}$).
All KINS reference values are traceable to national measurement standards.
Complete measured activity concentrations before and after solidification
are reported in Table~\ref{tab:rawdata}.

\begin{table}[htbp]
\footnotesize
\caption{Activity concentration data (Bq~kg$^{-1}$) transcribed from the
  laboratory summary sheet before and after LM-pectin--Ca$^{2+}$
  solidification.
  $A_{\text{KINS}}$: KINS proficiency-test certified reference value.
  $R$: paired recovery rate (Eq.~\ref{eq:recovery}).
  $\Delta$: paired change after solidification (Eq.~\ref{eq:change}).
  $B$: post-solidification bias relative to the KINS value
  (Eq.~\ref{eq:bias}).
  N/M: not measured.
  The original spectroscopy exports should accompany the public data
  package when available.}
\label{tab:rawdata}
\centering
\begin{tabular}{llcccccc}
\toprule
Radionuclide & Sample &
  $A_{\text{KINS}}$ & $A_{\text{before}}$ & $A_{\text{after}}$ &
  $R$ (\%) & $\Delta$ (\%) & $B$ (\%) \\
\midrule
\multirow{5}{*}{$^{137}$Cs}
  & S1 & 45.1 & 44.0  & 44.4  & 100.9 & $+$0.9 & $-$1.6 \\
  & S2 & 45.1 & 47.0  & 43.4  &  92.3 & $-$7.7 & $-$3.8 \\
  & S3 & 24.6 & 30.8  & 30.7  &  99.7 & $-$0.3 & $+$24.8 \\
  & S4 & 24.6 & 28.3  & 30.6  & 108.1 & $+$8.1 & $+$24.4 \\
  & S5 & 24.8 & 30.3  & 29.9  &  98.7 & $-$1.3 & $+$20.6 \\
\midrule
\multirow{5}{*}{$^{134}$Cs}
  & S1 & 45.1 & 44.0  & 40.2  &  91.4 & $-$8.6 & $-$10.9 \\
  & S2 & 45.1 & 47.0  & 47.3  & 100.6 & $+$0.6 & $+$4.9 \\
  & S3 & 24.6 & 30.8  & 28.4  &  92.2 & $-$7.8 & $+$15.4 \\
  & S4 & 24.6 & 28.3  & 20.1  &  71.0 & $-$29.0 & $-$18.3 \\
  & S5 & 24.8 & ---   & ---   & N/M   & N/M & N/M \\
\midrule
\multirow{5}{*}{$^{241}$Am}
  & S1 & 45.1 & 44.0  & 40.6  &  92.3 & $-$7.7 & $-$10.0 \\
  & S2 & 45.1 & 47.0  & 45.8  &  97.4 & $-$2.6 & $+$1.6 \\
  & S3 & 24.6 & 30.8  & 26.5  &  86.0 & $-$14.0 & $+$7.7 \\
  & S4 & 24.6 & 28.3  & 26.1  &  92.2 & $-$7.8 & $+$6.1 \\
  & S5 & 24.8 & 30.3  & 30.5  & 100.7 & $+$0.7 & $+$23.0 \\
\midrule
\multirow{5}{*}{$^{60}$Co}
  & S1 & 45.1 & 44.0  & 40.5  &  92.0 & $-$8.0 & $-$10.2 \\
  & S2 & 45.1 & 47.0  & 44.6  &  94.9 & $-$5.1 & $-$1.1 \\
  & S3 & 24.6 & ---   & ---   & N/M   & N/M & N/M \\
  & S4 & 24.6 & ---   & ---   & N/M   & N/M & N/M \\
  & S5 & 24.8 & 30.3  & 29.5  &  97.4 & $-$2.6 & $+$19.0 \\
\bottomrule
\end{tabular}
\end{table}

\subsection{Sample Preparation and Solidification Procedure}

\begin{figure}[ht]
  \centering
  \includegraphics[width=\linewidth]{Fig2_revised.png}
  \caption{Experimental procedure for LM-pectin--Ca$^{2+}$ sample solidification.
    (a) Cylindrical polyethylene measurement container
    (inner diameter 60~mm; fill height marker at 10~mm).
    (b) Gravimetric addition of CaCl$_2\cdot$2H$_2$O
    (0.0323~g; Sartorius analytical balance).
    (c) Stirring hotplate conditions (room temperature, 300~rpm).
    (d) Gradual addition of LM-pectin powder to the seawater sample.}
  \label{fig:procedure}
\end{figure}

The procedure is shown in Fig.~\ref{fig:procedure}.
We used a cylindrical polyethylene container (inner diameter 60~mm,
height 40~mm) as the counting geometry, with the fill height fixed
at 10~mm (28.26~g sample mass).
The 28.26~g sample mass corresponds to a 10~mm fill height in the
60~mm-diameter container and was selected as the fixed geometry used
for the ISOCS efficiency calculation.

For each sample: 28.26~g of KINS reference seawater was weighed into
the container (Fig.~\ref{fig:procedure}a); 0.0323~g of
CaCl$_2\cdot$2H$_2$O was weighed on a Sartorius balance
(readability 0.0001~g; Fig.~\ref{fig:procedure}b) and dissolved by
brief stirring; the container was placed on a digital stirring plate
(MSH-30D; Daihan Scientific, Wonju, Korea) set to 300~rpm at room
temperature (Fig.~\ref{fig:procedure}c); and LM-pectin powder was added
gradually over 2--3~min under continuous stirring
(Fig.~\ref{fig:procedure}d).

We checked gelation status every 10~min by briefly inverting the
container; gelation was considered complete when no flow was visible
within 5~s of inversion.
The gelled container was then placed directly on the HPGe detector
for measurement without any additional processing.

\subsection{Gelation Condition Optimisation}

Five pectin loadings were tested per 28.26~g seawater:
0.7, 1.4, 1.7, 2.8, and 3.5~g.
For each loading we recorded the time to gel formation and assessed
gel appearance by eye (transparency, surface smoothness, absence of
granulation or air inclusions) and by the inversion test.
We assigned each gel to one of five categories:
no gel, insufficient, optimal, over-gelled, or excessive.

\subsection{Gamma Spectrometry}

All measurements---before and after solidification---were made with
the same $p$-type coaxial HPGe detector (crystal diameter 70~mm,
height 80~mm; relative efficiency 30\%; energy resolution 1.85~keV
FWHM at 1332~keV; Canberra Industries, USA).
The detector was connected to a digital multi-channel analyser running
Genie-2000 spectroscopy software (Canberra).
Counting efficiency for the cylindrical container geometry was
calibrated using the ISOCS computational tool (Canberra).
Each measurement ran for 30000~s live time.

Four radionuclides were quantified by their principal gamma peaks:
$^{137}$Cs at 661.7~keV, $^{134}$Cs at 604.7~keV,
$^{241}$Am at 59.5~keV, and $^{60}$Co at 1173.2~keV.

\subsection{Recovery Rate Calculation}

Recovery $R$ was defined as:
\begin{equation}
  R\,(\%) = \frac{A_{\text{after}}}{A_{\text{before}}} \times 100
  \label{eq:recovery}
\end{equation}
where $A_{\text{before}}$ and $A_{\text{after}}$ (Bq~kg$^{-1}$) are
the activity concentrations measured before and after solidification.
The paired change after solidification was calculated as:
\begin{equation}
  \Delta\,(\%) = R - 100
  \label{eq:change}
\end{equation}
Agreement with the KINS certified reference value was reported separately
as the post-solidification bias:
\begin{equation}
  B\,(\%) = \frac{A_{\text{after}} - A_{\text{KINS}}}{A_{\text{KINS}}} \times 100
  \label{eq:bias}
\end{equation}
The separation of $\Delta$ and $B$ is important: $\Delta$ evaluates the
effect of the solidification step, whereas $B$ includes pre-existing
calibration and sample-specific differences relative to the certified value.


%% ══════════════════════════════════════════════════════════
\section{Results and Discussion}
%% ══════════════════════════════════════════════════════════

\subsection{Gelation Optimisation}

The gelation behaviour of LM-pectin in seawater varied systematically with
pectin amount (Fig.~\ref{fig:gelation}).
At 0.7~g per 28.26~g seawater, no gel formed within 180~min,
indicating insufficient cross-linking density at this pectin concentration.
At 1.4~g, a transparent, mechanically stable gel formed in 120~min,
exhibiting smooth surface topology and no phase separation;
this condition was designated as optimal.
Increasing the amount to 1.7~g reduced gelation time to 90~min
but produced a slightly opaque, firmer gel.
At 2.8~g and 3.5~g, gelation was effectively instantaneous ($<$5~min),
resulting in heterogeneous, granular gel structures with visible
air inclusions that would compromise the homogeneity required
for accurate gamma spectrometry.

\begin{figure}[htbp]
  \centering
  \includegraphics[width=0.7\linewidth]{Fig3.png}
  \caption{Gelation behaviour of LM-pectin in artificial seawater as a function
    of pectin amount (g per 28.26~g seawater).
    The bar chart shows gelation time (min) for each pectin loading;
    0.7~g produced no gel within 180~min, 1.4~g (optimal condition, green)
    formed a transparent, mechanically stable gel in 120~min,
    1.7~g gelled in 90~min with slightly reduced homogeneity,
    and 2.8~g and 3.5~g produced instantaneous, heterogeneous gel structures.
    Photographs below show the representative gel appearance for each condition
    (top view after inversion test).}
  \label{fig:gelation}
\end{figure}

\sloppy
The concentration dependence makes sense in terms of the egg-box
cross-linking mechanism~\cite{Braccini2001}.
Seawater Ca$^{2+}$ at $\sim$400~mg~L$^{-1}$ ($\sim$10~mmol~L$^{-1}$)
is enough to initiate junction-zone formation in LM-pectin, but only
above a threshold pectin concentration at which the chain density
is sufficient for a connected network to form.
At 1.4~g per 28.26~g ($\approx$4.7~wt\%), we appear to be just above
that threshold.
At 0.7~g ($\approx$2.5~wt\%), we are below it---the Ca$^{2+}$ is
present but the pectin chains are too dilute to form a percolating
network within 180~min.
At 2.8~g and above, the opposite problem occurs: pectin chains are
dense enough to saturate available Ca$^{2+}$ binding sites almost
instantaneously, trapping the polymer in a heterogeneous state before
the chains can arrange themselves into a uniform gel.
The 1.4~g condition avoids both problems and was selected for all
subsequent experiments.

\subsection{Cs-137 Recovery: Individual Sample Analysis}

Fig.~\ref{fig:recovery}a shows $^{137}$Cs concentrations measured before
and after solidification for all five samples (three KINS rounds:
2019, 2022, 2024).
In every case, the after-solidification value stayed within a few
percent of the before-solidification value across reference
concentrations from 24.6 to 45.1~Bq~kg$^{-1}$.
Mean paired recovery was 99.9~$\pm$~5.6\%; the largest absolute paired
change was 8.1\% (S4).

The near-complete recovery of $^{137}$Cs is not surprising.
Cs$^+$ is a monovalent cation that stays dissolved in the gel pore
water; it has no affinity for the carboxylate groups on the pectin
backbone, which selectively bind divalent cations like Ca$^{2+}$
and Sr$^{2+}$~\cite{Braccini2001}.
The seawater Ca$^{2+}$ at $\sim$400~mg~L$^{-1}$ already saturates most
available binding sites, so trace-level $^{137}$Cs$^+$ is effectively
excluded from the cross-link structure and remains evenly distributed
in the water phase of the gel.

\subsection{Multi-Radionuclide Recovery}

\begin{figure}[htbp]
  \centering
  \includegraphics[width=\linewidth]{Fig4_revised.png}
  \caption{$^{137}$Cs recovery performance and multi-radionuclide recovery
    summary.
    (a) $^{137}$Cs activity concentrations before (open circles) and after
    (filled circles) LM-pectin--Ca$^{2+}$ solidification for five samples
    (S1--S5) spanning three KINS proficiency-test rounds.
    (b) Mean recovery rates ($\pm$1~SD) for $^{137}$Cs ($n=5$),
    $^{134}$Cs ($n=4$), $^{241}$Am ($n=5$), and $^{60}$Co ($n=3$)
    after LM-pectin--Ca$^{2+}$ solidification.
    The dashed line indicates 100\% recovery and dotted lines mark the
    80--120\% paired-recovery screening range.}
  \label{fig:recovery}
\end{figure}

\begin{table}[htbp]
\caption{Paired recovery rates of four radionuclides
measured after LM-pectin--Ca$^{2+}$ solidification.
Individual sample-level data are provided in Table~\ref{tab:rawdata}.}
\label{tab:recovery}
\centering
\begin{tabular}{lccccc}
\toprule
Radionuclide & $E_\gamma$ (keV) & $n$ &
  Mean recovery (\%) & Minimum (\%) & Maximum (\%) \\
\midrule
$^{137}$Cs & 661.7   & 5 & $99.9 \pm 5.6$  & 92.3 & 108.1 \\
$^{134}$Cs & 604.7   & 4 & $88.8 \pm 12.6$ & 71.0 & 100.6 \\
$^{241}$Am & 59.5    & 5 & $93.7 \pm 5.6$  & 86.0 & 100.7 \\
$^{60}$Co  & 1173.2  & 3 & $94.8 \pm 2.7$  & 92.0 & 97.4 \\
\bottomrule
\multicolumn{6}{l}{\footnotesize Values are mean $\pm$ 1 standard deviation
  (sample SD, $ddof=1$).}
\end{tabular}
\end{table}

Fig.~\ref{fig:recovery}b and Table~\ref{tab:recovery} summarise
mean recovery for all four nuclides.
Sixteen of the 17 paired values were within the 80--120\% screening
range. Sample-level data are in Table~\ref{tab:rawdata}.

$^{134}$Cs recovered at 88.8~$\pm$~12.6\%, with one low paired result
(S4, 71.0\%). The current dataset does not establish the cause of this
deviation. The result is retained in the analysis and treated as a
limitation rather than excluded as an outlier.

$^{60}$Co recovered at 94.8~$\pm$~2.7\%, with the narrowest
range of the four nuclides ($n=3$).
Only three of the five KINS rounds included $^{60}$Co at detectable
concentrations, so $n=3$ here.
All three paired values were within 92.0--97.4\%, supporting further
evaluation of the gel protocol at high energies.

The 93.7\% recovery for $^{241}$Am at 59.5~keV is worth a closer look.
At this energy, photons are partially absorbed by the sample before
reaching the detector.
Water (and the gel, which is $>$95\% water by mass) has a mass
attenuation coefficient of roughly 0.20~cm$^2$~g$^{-1}$ at
59.5~keV~\cite{Hubbell2004}, which translates to a 2--5\%
self-absorption correction for a 10~mm fill height.
The observed 93.7\% recovery is compatible with a modest low-energy
attenuation effect, although the present dataset cannot separate
self-absorption from other sample-preparation effects.
Users measuring $^{241}$Am should validate an energy-dependent efficiency
correction for the selected fill height.

\subsection{Paired Variation Across the Energy Range}

\begin{figure}[ht]
  \centering
  \includegraphics[width=\linewidth]{Fig5_revised.png}
  \caption{(a) Heatmap of individual paired changes $\Delta$ (\%)
    for each radionuclide--sample combination
    (Eq.~\ref{eq:change}).
    N/M: not measured for that sample--nuclide pair.
    (b) Mean $^{60}$Co paired recovery with 95\% confidence interval
    ($n=3$, $ddof=1$). The shaded area indicates the 80--120\%
    paired-recovery screening range.}
  \label{fig:heatmap}
\end{figure}

Fig.~\ref{fig:heatmap}a maps the paired change caused by the
solidification step. The low-energy $^{241}$Am measurements remained
within $-14.0$ to $+0.7$\%, whereas the most pronounced deviation was
the $^{134}$Cs S4 result ($-29.0$\%). With the present sample size,
the dataset is not sufficient for a formal energy-dependence claim.

For $^{60}$Co, the mean paired recovery was 94.8\% and the 95\%
confidence interval was 88.2--101.4\% (Fig.~\ref{fig:heatmap}b).
A descriptive sensitivity analysis based on the observed SD indicates
that the $n=3$ dataset had 80\% power to detect an absolute shift of
approximately 8.7 percentage points in a two-sided one-sample
$t$-test ($\alpha=0.05$). The analysis is provided in the Supplementary
Information with the raw values and code. It does not replace future
prospective replication.

\subsection{Field Applicability}

The central problem this method addresses is straightforward: liquid
seawater in a counting container sloshes with every roll and pitch of
the vessel, so the distance from the sample surface to the HPGe detector
crystal changes continuously during measurement.
Because HPGe counting efficiency falls off with source-to-detector
distance, even moderate swell introduces counting efficiency errors
that are difficult to correct retrospectively without knowing the exact
geometry at every moment of the measurement.
Once the seawater is converted to a gel, the sample surface is fixed and
the geometry is defined. This should reduce sensitivity to vessel motion,
although direct testing aboard a vessel or on a motion platform remains
necessary.

A secondary advantage is operational safety and handling convenience.
A gel is less likely to spill than a liquid sample.
In a confined shipboard instrument room, a spilled liquid sample
containing elevated $^{137}$Cs or other radionuclides would require
immediate decontamination and could damage the HPGe detector; a gel
remains in the container even if the container tips.
The prepared gels showed no visible syneresis during the measurement
window. Longer stability tests remain necessary for workflows that
require storage or transport.

From a logistics perspective, the reagents needed are minimal.
LM-pectin is a food-grade powder, and both reagents can be stored at
room temperature.
The procedure needs only a stirring hotplate and an analytical balance;
no fume hood, no pre-evaporation step, and no specialised chemistry skills
are required.
Gelation is complete within 2~h under the selected condition.

In comparison with other solidification approaches, LM-pectin offers
some additional practical points.
Cement-based matrices are designed for durable immobilisation and are
irreversible once set~\cite{Rahman2007}.
Epoxy fixation is also a durable approach rather than a reversible
sample-preparation workflow~\cite{Costa2021}.
Some synthetic hydrogel routes require reagents and preparation steps
that are less convenient for simple field handling~\cite{Du2020}.
The LM-pectin--Ca$^{2+}$ gel, by contrast, uses a food-grade polymer
and a simple aqueous preparation protocol~\cite{May1990}.

\subsection{Comparison with Alginate-Based Gelation}

Sodium alginate gels in seawater by the same egg-box mechanism as
LM-pectin~\cite{Grant1973, Makarova2023}, so it is worth asking
whether alginate would work equally well as a solidification matrix.

The literature on alginate in radioactivity research points in a
different direction: alginate has been studied mainly as a material
for \emph{removing} radionuclides from water, not for preserving
their concentration.
Prussian blue--alginate aerogels, for example, remove $\sim$20~mg~g$^{-1}$
of Cs$^+$ and Sr$^{2+}$ from seawater~\cite{Eun2021} by exploiting
strong cation binding to the polymer backbone.
Pectin-based sorbents have been explored for the same purpose~\cite{BokBadura2020}.
Strong cation affinity is useful for decontamination materials but may
redistribute radionuclides within a measurement matrix. The 99.9\%
mean paired recovery for $^{137}$Cs indicates that substantial net loss
did not occur under the LM-pectin preparation conditions used here.
Spatial distribution within the gel was not measured directly.

Whether alginate under the same conditions (seawater Ca$^{2+}$ as the
cross-linker) would show similarly low Cs binding has not been tested
directly.
Direct comparison with alginate was outside the scope of the present
study and would be useful in future protocol-development work.
More broadly, radionuclide coordination is ligand- and ion-dependent,
as shown in computational chelator-selection studies~\cite{Kim2026}.

\subsection{Limitations and Future Directions}

Four limitations of this study are worth stating clearly.

First, all experiments used ASTM D1141 artificial seawater, not
real coastal or open-ocean seawater.
Real seawater carries dissolved organic matter, particulates,
and microbial material that are absent from the artificial standard.
Whether these components affect gelation kinetics or cause
radionuclide redistribution within the gel matrix is not known.
The pH of real seawater (7.8--8.3) is slightly different from
that of the artificial standard too, and pH affects the degree of
pectin deprotonation and therefore the rate of Ca$^{2+}$ binding.
Validation with real contaminated samples from active coastal monitoring
stations is the most important next step~\cite{Choi2023}.

Second, motion tolerance was not tested aboard a vessel or with a
laboratory motion simulator. The field benefit is therefore a practical
hypothesis supported by the fixed gel geometry, not a completed
shipboard validation.

Third, the reference-bias analysis identified four post-solidification
measurements outside $\pm20$\% of the KINS certified value
(Table~\ref{tab:rawdata}). Because discrepancies were also present
before solidification, paired recovery is the primary endpoint for this
protocol study. Absolute calibration performance should be re-evaluated
from the original spectroscopy exports before operational deployment.

Fourth, the replicate count for $^{60}$Co was only $n=3$, because
only three of the five KINS reference rounds included this nuclide
at detectable concentrations.
The confidence interval and descriptive sensitivity analysis are reported
to make this limitation explicit.

Looking ahead, we see three useful extensions of this work:
extending the nuclide panel to naturally occurring $^{40}$K and
$^{7}$Be, which are of independent interest for marine atmospheric
deposition studies~\cite{Bae2025}; developing Monte Carlo
(MCNP/GEANT4) efficiency correction models specific to the gel
geometry and water-equivalent matrix to tighten accuracy at low
gamma energies~\cite{Corbacho2022}; and testing the method
aboard an actual vessel during a monitoring mission to confirm
that the geometry stability benefit translates to real measurement
conditions.


%% ══════════════════════════════════════════════════════════
\section{Conclusions}
%% ══════════════════════════════════════════════════════════

Solidifying seawater samples with LM-pectin and controlled Ca$^{2+}$
supplementation addresses a practical sample-handling problem:
liquid samples can change geometry and spill during mobile-platform
measurement. The gel fixes the sample shape and provides a basis for
future shipboard validation.

At the optimised condition---1.4~g LM-pectin per 28.26~g seawater,
stirred at 300~rpm for 120~min at room temperature---radionuclide
mean paired recovery was
99.9\%, 88.8\%, 93.7\%, and 94.8\% for $^{137}$Cs, $^{134}$Cs,
$^{241}$Am, and $^{60}$Co, respectively, across 59.5--1173.2~keV.
Sixteen of 17 paired values remained within 80--120\% recovery.
The $^{134}$Cs S4 result (71.0\%) and reference-bias discrepancies are
reported transparently as limitations (Table~\ref{tab:rawdata};
Figs.~\ref{fig:gelation}--\ref{fig:heatmap}).
No pre-evaporation or specialised reagent is needed.

The present experiments used ASTM artificial seawater and did not
include a vessel-motion test. Confirming performance with real seawater,
evaluating motion tolerance, extending the nuclide panel to $^{40}$K
and $^{7}$Be~\cite{Bae2025}, and developing Monte Carlo efficiency
corrections are the practical next steps toward deployment.


\nolinenumbers
%% ══════════════════════════════════════════════════════════
\section*{CRediT Author Contribution Statement}
%% ══════════════════════════════════════════════════════════
\textbf{Yejin Lee:} Conceptualisation, Methodology, Investigation,
Formal Analysis, Writing -- Original Draft.
\textbf{Hyunsung Kim:} Methodology, Formal Analysis, Validation,
Writing -- Review \& Editing.
\textbf{Myeongsik Shin:} Investigation, Resources, Writing -- Review \& Editing.
\textbf{Gyoung Jun Choi:} Investigation, Resources.
\textbf{Jun Yoo:} Investigation, Data Curation.
\textbf{Ku Kang:} Conceptualisation, Supervision, Funding Acquisition,
Writing -- Review \& Editing.

\section*{Declaration of Competing Interest}
The authors declare that they have no known competing financial
interests or personal relationships that could have appeared to
influence the work reported in this paper.

\section*{Data Availability Statement}
The data table, derived metrics, figure-reproduction code, and
$^{60}$Co sensitivity analysis are available at:
\url{[DATA_REPOSITORY_URL_TO_BE_ADDED_BEFORE_SUBMISSION]}.
The spectroscopy export files should be deposited in the same repository
before submission.

\section*{Supplementary Information}
Supplementary Information includes the calculation definitions,
the $^{60}$Co sensitivity analysis, and a data-file inventory.

\section*{Acknowledgements}
The KINS proficiency test reference standards (2019, 2022, 2024 rounds)
are gratefully acknowledged.



%% ══════════════════════════════════════════════════════════
%% REFERENCES
%% ══════════════════════════════════════════════════════════
\begin{thebibliography}{10}
\expandafter\ifx\csname url\endcsname\relax
  \def\url#1{\texttt{#1}}\fi
\expandafter\ifx\csname urlprefix\endcsname\relax\def\urlprefix{URL }\fi
\expandafter\ifx\csname href\endcsname\relax
  \def\href#1#2{#2} \def\path#1{#1}\fi

\bibitem{Buesseler2012}
K.~O. Buesseler, S.~R. Jayne, N.~S. Fisher, I.~I. Rypina, H.~Baumann, Z.~Baumann, C.~F. Breier, E.~M. Douglass, J.~George, A.~M. Macdonald, H.~Miyamoto, J.~Nishikawa, S.~M. Pike, S.~Yoshida, Fukushima-derived radionuclides in the ocean and biota off {Japan}, Proceedings of the National Academy of Sciences 109~(16) (2012) 5984--5988.
\newblock \href {https://doi.org/10.1073/pnas.1120794109} {\path{doi:10.1073/pnas.1120794109}}.

\bibitem{Lu2021}
Y.~Lu, J.~Yuan, D.~Du, B.~Sun, X.~Yi, Monitoring long-term ecological impacts from release of fukushima radiation water into ocean, Geography and Sustainability 2~(2) (2021) 95--98.

\bibitem{IAEA2005}
{IAEA}, Worldwide Marine Radioactivity Studies ({WOMARS}): Radionuclide Levels in Oceans and Seas, International Atomic Energy Agency, Vienna, Austria, 2005, iAEA-TECDOC-1429.

\bibitem{Abbasi2021}
A.~Abbasi, H.~M.~H. Zakaly, A.~Badawi, The anthropogenic radiotoxic element of $^{137}${Cs} accumulate to biota in the {Mediterranean Sea}, Marine Pollution Bulletin 164 (2021) 112043.
\newblock \href {https://doi.org/10.1016/j.marpolbul.2021.112043} {\path{doi:10.1016/j.marpolbul.2021.112043}}.

\bibitem{Uddin2020}
S.~Uddin, S.~W. Fowler, M.~Behbehani, A.~N. Al-Ghadban, P.~W. Swarzenski, N.~Al-Awadhi, A review of radioactivity in the {Gulf} region, Marine Pollution Bulletin 159 (2020) 111481.
\newblock \href {https://doi.org/10.1016/j.marpolbul.2020.111481} {\path{doi:10.1016/j.marpolbul.2020.111481}}.

\bibitem{Jerome2014}
S.~M. Jerome, M.~J. Woods, S.~E.~M. Lucas, L.~Johansson, J.~F. Gravel, The role of naturally occurring radionuclides in the standardisation of low level radioactivity measurement, Applied Radiation and Isotopes 87 (2014) 302--308.
\newblock \href {https://doi.org/10.1016/j.apradiso.2013.11.079} {\path{doi:10.1016/j.apradiso.2013.11.079}}.

\bibitem{Debertin1988}
K.~Debertin, R.~G. Helmer, Gamma- and {X}-ray Spectrometry with Semiconductor Detectors, North-Holland, Amsterdam, 1988.


\bibitem{Corbacho2022}
J.~A. Corbacho, A.~Baeza, In situ gamma spectrometry using a portable {HPGe} detector: radiological characterisation and environmental surveillance around an operating nuclear power plant, Journal of Radiological Protection 42~(1) (2022) 011501.
\newblock \href {https://doi.org/10.1088/1361-6498/ac4116} {\path{doi:10.1088/1361-6498/ac4116}}.

\bibitem{Choi2024}
Y.~Choi, Y.-Y. Ji, S.~Joung, B.-I. Min, Development of a portable underwater in-situ radiation detector ({MARK-U4}) for marine radioactivity monitoring, Journal of Coastal Research 116 (2024) 176--180.
\newblock \href {https://doi.org/10.2112/JCR-SI116-036.1} {\path{doi:10.2112/JCR-SI116-036.1}}.


\bibitem{Barrera2017}
M.~Barrera, M.~Casas-Ruiz, J.~J. Alonso, J.~Vidal, Precise determination of hpge detector efficiency for gamma spectrometry measurements of environmental samples with variable geometry and density, Nukleonika 62~(1) (2017) 47--59.

\bibitem{Choi2023}
Y.~Choi, Y.-Y. Ji, S.~Joung, Experimental counting efficiency calibration of in situ marine radioactivity monitoring systems in {Busan} ({Korea}) based on seawater salinity, Journal of Coastal Research 39 (2023) 377--383.
\newblock \href {https://doi.org/10.2112/JCOASTRES-D-22TM-00013.1} {\path{doi:10.2112/JCOASTRES-D-22TM-00013.1}}.


\bibitem{Rahman2007}
R.~A. Rahman, A.~Zaki, A.~El-Kamash, Modeling the long-term leaching behavior of 137cs, 60co, and 152,154 eu radionuclides from cement--clay matrices, Journal of Hazardous Materials 145~(3) (2007) 372--380.


\bibitem{Du2020}
H.~Du, S.~Shi, W.~Liu, H.~Teng, M.~Piao, Processing and modification of hydrogel and its application in emerging contaminant adsorption and in catalyst immobilization: a review, Environmental Science and Pollution Research 27~(12) (2020) 12967--12994.

\bibitem{Braccini2001}
I.~Braccini, S.~Perez, Molecular basis of {Ca}$^{2+}$-induced gelation in alginates and pectins: the egg-box model revisited, Biomacromolecules 2~(4) (2001) 1089--1096.
\newblock \href {https://doi.org/10.1021/bm010008g} {\path{doi:10.1021/bm010008g}}.

\bibitem{Grant1973}
G.~T. Grant, E.~R. Morris, D.~A. Rees, P.~J.~C. Smith, D.~Thom, Biological interactions between polysaccharides and divalent cations: the egg-box model, FEBS Letters 32~(1) (1973) 195--198.
\newblock \href {https://doi.org/10.1016/0014-5793(73)80770-7} {\path{doi:10.1016/0014-5793(73)80770-7}}.

\bibitem{Millero2008}
F.~J. Millero, R.~Feistel, D.~G. Wright, T.~J. McDougall, The composition of {Standard Seawater} and the definition of the {Reference-Composition Salinity Scale}, Deep-Sea Research Part I 55~(1) (2008) 50--72.
\newblock \href {https://doi.org/10.1016/j.dsr.2007.10.001} {\path{doi:10.1016/j.dsr.2007.10.001}}.

\bibitem{May1990}
C.~D. May, Industrial pectins: Sources, production and applications, Carbohydrate Polymers 12 (1990) 79--99.
\newblock \href {https://doi.org/10.1016/0144-8617(90)90105-2} {\path{doi:10.1016/0144-8617(90)90105-2}}.

\bibitem{Hubbell2004}
J.~H. Hubbell, S.~M. Seltzer, Tables of {X}-Ray Mass Attenuation Coefficients and Mass Energy-Absorption Coefficients (version 1.4), National Institute of Standards and Technology, Gaithersburg, MD, 2004, available at: \url{https://physics.nist.gov/xaamdi}.

\bibitem{Makarova2023}
A.~O. Makarova, S.~R. Derkach, T.~Khair, M.~A. Kazantseva, Y.~F. Zuev, O.~S. Zueva, Ion-induced polysaccharide gelation: Peculiarities of alginate egg-box association with different divalent cations, Polymers 15~(5) (2023) 1243.
\newblock \href {https://doi.org/10.3390/polym15051243} {\path{doi:10.3390/polym15051243}}.

\bibitem{Eun2021}
S.~Eun, J.~Ryu, H.~Kim, H.-J. Hong, S.~Kim, Simultaneous removal of radioactive cesium and strontium from seawater using a highly efficient {Prussian blue}-embedded alginate aerogel, Journal of Environmental Management 297 (2021) 113389.
\newblock \href {https://doi.org/10.1016/j.jenvman.2021.113389} {\path{doi:10.1016/j.jenvman.2021.113389}}.

\bibitem{BokBadura2020}
J.~Bok-Badura, A.~Jak\'{o}bik-Kolon, A.~Kazek-K\k{e}sik, K.~Karo\'{n}, Hybrid pectin-based sorbents for cesium ion removal, Materials 13~(9) (2020) 2160.
\newblock \href {https://doi.org/10.3390/ma13092160} {\path{doi:10.3390/ma13092160}}.

\bibitem{Bae2025}
S.~Bae, Y.~Lee, S.~Lee, K.-D. Byun, K.~Kang, Environmental drivers of be-7 and k-40 deposition in south korea: A regional analysis of meteorological and geological influences, J. Environ. Radioact. 288 (2025) 107725.
\newblock \href {https://doi.org/10.1016/j.jenvrad.2025.107725} {\path{doi:10.1016/j.jenvrad.2025.107725}}.

\bibitem{Kim2026}
J.~Kim, K.~Kang, K.~Shin, M.~K. Schultz, W.~B. Lee, D.~Lee,
Rational selection of chelators for theranostic radionuclides:
Insights from ab initio computational modeling,
Molecular Pharmaceutics 23 (2026) 2779--2792.
\newblock \href {https://doi.org/10.1021/acs.molpharmaceut.5c01955}
{\path{doi:10.1021/acs.molpharmaceut.5c01955}}.

\bibitem{Costa2021}
O.~L. da~Costa, I.~S.~M. Freitas, F.~A. Genezini, D.~S. Moreira, G.~S. Zahn,
Gamma spectrometry of iodine-125 produced in {IEA-R1} nuclear reactor,
using {HPGe} detector and fixation into epoxy matrix disc,
Applied Radiation and Isotopes 169 (2021) 109568.
\newblock \href {https://doi.org/10.1016/j.apradiso.2020.109568} {\path{doi:10.1016/j.apradiso.2020.109568}}.

\end{thebibliography}


\end{document}
